\honest{Sympathetic Minds}{Eliezer Yudkowsky}{2009}

Mirror neurons are active both when you perform an action and when you watch someone else perform it. They have been recorded in primates, and imaging has found consistent evidence in humans. \nb{That was an accurate summary in 2009. Recordings from single human neurons followed in 2010.}

As I wrote before, brains are so complex that if we tried to understand them from scratch, as we do gravity or a car, we could not invent good hypotheses in a lifetime. The only way to hit on an ``Aha!'' about another mind is to force something very similar, your own brain, to behave similarly. That is what I call ``empathy''. ``Sympathy'' goes further: to smile when someone else smiles and to hurt when they hurt, moving from prediction into reinforcement. \nb{The words are used in the author's own senses. In psychology, sharing another's feeling is usually called empathy, and sympathy means concern for the other.}

Why would natural selection make anything so kind? It might have started, maybe, with a mother's love for her children or a sibling's: smiling when they smile and wincing when they wince leads you to help in many ways, and in the ancestral environment what relatives want probably bore on their reproductive success. Evolution does not take ``the elaborate correct optimal path'' of an abstract desire for your relatives to get what they want. The mirroring machinery was already there, so sympathy was a short step. An abstract desire would have needed a whole definition of ``wanting''; evolution falls up the fitness landscape like water flowing downhill.

Then came reciprocity with allies, Tit for Tat elaborated for reputations. I think the most formidable person is more often the one with the most friends. A general urge to smile when friends smile wins more of them than a specific urge like the vampire bat's to share blood, and an organism that winces when friends wince will avoid angering them. Of course you also want to kill designated enemies without a qualm. These are humans. \nb{Frans de Waal had argued much the same in 2008, proposing shared emotional states as the mechanism of altruism toward those in distress, one that fits both kin selection and reciprocity.}

I am not sure, but based on my ``admittedly limited knowledge of anthropology'' sympathy looks ``on'' by default in humans: some cultures help strangers and some eat them, and the question is which needs the explicit rule. It is painful to be a bystander in a war whose two sides have switched off sympathy for each other while yours is on for both, so that you wince at a dead child whatever the caption.

So sympathy, a strange and deep implementation of reciprocity and help, tangles minds together not by a term in the utility function for another's desire but ``by the simpler and yet far more consequential path of mirror neurons.'' \nb{Here the post drops its hedges. The link between mirror neurons and empathy was disputed in 2009, and a 2015 review of empathy research still questioned it.}

An AI need not work like this. Empathy is one way of predicting minds, not the only one. The human brain cannot quickly rewire visual cortex for hearing in a dark room, but an AI can shift resources at once and swap programs to disk. So it need not force its own mind into a similar state, with all the risk and mess of mixing data with its own state. It can create a separate mind-instance, or model the other mind directly, as a hypothesis like any other. You don't need to become your shoes to understand your shoes.

A paperclip maximizer that needs humans would model their reactions and manage them, and might smile to manipulate them, but would feel nothing. The word ``happiness'' is an idiom of reinforcement learning, not of expected utility maximization, and the paperclipper simply chooses the action with the most expected paperclips. Hating an enemy is not the right mode either.

To imagine how it sees you, picture yourself as a machine with levers, not human-shaped, say a woodsaw: some levers produce coins, some fire a bullet, and the levers must be pulled in the right order. To understand unsympathetic optimization, study natural selection, which does not anaesthetize dying creatures, since the anaesthetic would serve no reproductive purpose either.

That is why I put sympathy ahead of even boredom on my list of what would make aliens sympathetic. Mirror neurons, having happened once, could happen again. Unsympathetic aliens might trade with us, keep treaties and even cooperate in the Prisoner's Dilemma, but they would never be friends. They would see us only as means, never shed a tear for us or smile at our joys, give their own kind no different consideration, and feel nothing missing. They would be ``varelse,'' not ``ramen'': aliens we cannot relate to on any personal level, with no point in trying. \nb{The terms are from Orson Scott Card's \textsc{Speaker for the Dead}.}
